% TOP_PROBLEM: {"problemId": "problem.montgomery-s-large-value-conjecture-for-dirichlet-polynomials", "canonicalTitle": "Montgomery's Large Value Conjecture for Dirichlet Polynomials", "releaseRank": 295, "primaryDomain": "number_theory_arithmetic_geometry", "primaryDomainLabel": "Number theory & arithmetic geometry", "displayStatus": "open_with_solved_subcases", "releaseStatus": "open_with_solved_subcases"}
% !TeX program = lualatex
% Definition and sources only; this file is not an LLM solution attempt.
\documentclass[11pt]{article}
\usepackage[margin=25mm]{geometry}
\usepackage{fontspec}
\setmainfont{Latin Modern Roman}
\usepackage{amsmath,amssymb,mathtools}
\usepackage{unicode-math}
\setmathfont{Latin Modern Math}
\usepackage{xurl,hyperref}
\hypersetup{colorlinks=true,urlcolor=blue}
\setcounter{secnumdepth}{0}
\setlength{\parindent}{0pt}
\setlength{\parskip}{5pt}
\setlength{\emergencystretch}{3em}
\begin{document}
\section*{\texorpdfstring{Montgomery's Large Value Conjecture for Dirichlet Polynomials}{Montgomery's Large Value Conjecture for Dirichlet Polynomials}}\label{295}
\subsection{Definitions and mathematical statement}
Let $D(t)=\sum_{N<n\le2N}b_n n^{it}$ with $|b_n|\le1$, and let $W\subset[0,T]$ be one-separated: distinct elements differ by at least one. For fixed $\sigma>1/2$, assume $|D(t)|>N^\sigma$ on $W$. The target is
\[
 |W|\le T^{o(1)}N^{2-2\sigma},\qquad N\le T\le N^{O(1)}.
\]
Explicitly, with a fixed exponent range $T\le N^A$, the subpower notation asks for every $\eta>0$ a bound
$|W|\le C_{\sigma,A,\eta}T^\eta N^{2-2\sigma}$ uniformly in the coefficients and the separated set. An exponent range growing with $N$ is not the fixed-polynomial-range convention. The statement counts large values, not only an integral mean value.
\par\smallskip\textit{Formulation and concept references:} \href{https://arxiv.org/html/2405.20552}{[S1]}.

\subsection{Short English statement}
Can a bounded-coefficient Dirichlet polynomial take very large values at only the conjectured number of well-separated times in every polynomial-length interval?
\subsection{Sources}
\begin{itemize}
\item[S1]Larry Guth and James Maynard, New large value estimates for Dirichlet polynomials, Annals of Mathematics 203 (2026), 623-675, Conjecture 1.5.
\url{https://arxiv.org/html/2405.20552}.
\textit{Formulation source.}
\end{itemize}
\end{document}
